\section{Complete statistical criteria}\label{sec:criteria}
The cost of \emph{selecting} a smoothing level is not the cost of one fit. For fixed numeric $\mu$ and finite $\lambda>0$, write $e_\lambda=y-\widehat\tau_{\lambda,\mu}$, $\RSS_\lambda=e_\lambda^\top e_\lambda$, $h_i=(H_\lambda)_{ii}$, and $\nu=\tr H_\lambda$. The exact classical leave-one-out and generalized cross-validation scores are
\begin{equation}\label{eq:cvgcv}
\CV(\lambda)=\frac1n\sum_{i=1}^n
\left(\frac{e_{\lambda,i}}{1-h_i}\right)^2,\qquad
\GCV(\lambda)=\frac{\RSS_\lambda/n}{(1-\nu/n)^2}.
\end{equation}
These are retrospective smooth-trend criteria, not chronological forecasting-validation scores. The fixed-$\mu$ leverage diagonal is available spectrally as
\begin{equation}\label{eq:leverage}
h_i=1-\sum_{j=1}^r V_{ij}^2 w_j(\lambda).
\end{equation}
The unpenalized polynomial projection is implicitly included by the leading one in \eqref{eq:leverage}. Omitting it would give incorrect CV denominators.

For comparability with the conventions in \citet{cortes2017}, we also consider the residual-based expressions
\begin{equation}\label{eq:ic}
\begin{aligned}
\mathrm{AICc}_{\RSS}(\lambda)
 &=\log(\RSS_\lambda/n)+\frac{2\nu+1}{n-\nu-2},\\
\mathrm{BIC}_{\RSS}(\lambda)
 &=\log(\RSS_\lambda/n)+\frac{\nu\log n}{n}.
\end{aligned}
\end{equation}
The AICc value is defined only when its denominator is admissible, and the names here refer to precisely the stated residual-based conventions, not to every information criterion with the same name in other statistical models.

\begin{proposition}[Residual-BIC interpolation pathology]
\label{prop:bic}
For fixed $\mu$ suppose $Q_d(y-p_\mu)\ne0$. Then the expression $\mathrm{BIC}_{\RSS}(\lambda)$ in \eqref{eq:ic} tends to $-\infty$ as $\lambda\downarrow0$, and hence as $s=S_d(\lambda)\downarrow0$.
\end{proposition}
\begin{proof}
Since $H_\lambda=I-\lambda Q_d+O(\lambda^2)$,
$e_\lambda=(I-H_\lambda)(y-p_\mu)=
\lambda Q_d(y-p_\mu)+O(\lambda^2)$.
Thus $\RSS_\lambda=\lambda^2\|Q_d(y-p_\mu)\|^2+O(\lambda^3)$, while $\nu\to n$. The logarithmic residual term equals $2\log\lambda+O(1)$ and dominates the finite degrees-of-freedom penalty. The smoothness limit follows from Theorem~\ref{thm:bijection}.
\end{proof}
A strictly positive imposed grid minimum may be a disclosed sensitivity restriction, but is not a repaired global optimum for this expression. Likewise, CV/GCV near exact interpolation require careful limiting conventions rather than evaluating an indeterminate ratio at $s=0$.

\subsection{Gaussian reduced-coordinate likelihood}
To distinguish statistical selection from penalized training residuals, suppose in addition that
\begin{equation}\label{eq:gaussian}
y=\tau+\varepsilon,\quad
\varepsilon\sim N(0,\sigma_\varepsilon^2I_n),\quad
D_d\tau-\mu\mathbf1_r\sim N(0,\sigma_u^2I_r),
\quad \lambda=\sigma_\varepsilon^2/\sigma_u^2,
\end{equation}
with independent components and a fixed numeric $\mu$. The $r$ penalized SVD coordinates $z=V^\top(y-p_\mu)$ have marginal variances $\sigma_\varepsilon^2c_j(\lambda)$, where $c_j=1+1/(\lambda\delta_j)$. Profiling $\sigma_\varepsilon^2$ gives the criterion
\begin{equation}\label{eq:ml}
\mathcal L(\lambda)=
r\log\left(\frac{1}{r}\sum_{j=1}^r
\frac{z_j^2}{c_j(\lambda)}\right)
+\sum_{j=1}^r\log c_j(\lambda),
\end{equation}
up to a constant independent of $\lambda$. It is a reduced-coordinate variance-component likelihood, not the result of minimizing the PLS training objective with respect to $\lambda$.

\begin{proposition}[Banded likelihood sufficiency]\label{prop:det}
Let $\lambda>0$, $B_d=D_dD_d^\top$, and let $\widehat\tau$ solve \eqref{eq:rhs}. Set
$J_\lambda=\|y-\widehat\tau\|^2+
\lambda\|D_d\widehat\tau-\mu\mathbf1_r\|^2$.
Then
\begin{equation}\label{eq:quadraticidentity}
J_\lambda=\sum_{j=1}^r\frac{z_j^2}{c_j(\lambda)},
\end{equation}
and
\begin{equation}\label{eq:detidentity}
\sum_{j=1}^r\log c_j(\lambda)
=\log\det A_\lambda-r\log\lambda-\log\det B_d.
\end{equation}
Consequently a banded factorization of $A_\lambda$ and a precomputed banded factorization of $B_d$ suffice to evaluate \eqref{eq:ml} without a dense inverse.
\end{proposition}
\begin{proof}
In the positive spectral directions the minimal value of the quadratic PLS objective is $\sum_j w_j z_j^2$. Since $w_j=c_j^{-1}$ this proves \eqref{eq:quadraticidentity}; the nullspace part of the objective vanishes at the optimum. The positive eigenvalues of $Q_d$ and $B_d$ coincide. Thus
$\log\det A_\lambda=\sum_j\log(1+\lambda\delta_j)$ and $\log\det B_d=\sum_j\log\delta_j$. Subtracting $r\log\lambda$ proves \eqref{eq:detidentity}. Both determinants admit banded Cholesky log-diagonal evaluations.
\end{proof}

The same likelihood formula must not be called joint unknown-$\mu$ maximum likelihood after merely substituting $\widehat\mu=\operatorname{mean}(D_dy)$. Such a plug-in policy changes the statistical estimation problem. The direct PLS fit, likelihood fit, and genuinely foldwise-refitted leave-one-out selector are distinct objects.
